\documentclass{article}
\usepackage{amsmath,amssymb}
\usepackage{xurl}
\usepackage[hidelinks]{hyperref}
\setlength{\emergencystretch}{2em}
% Shared commands for the notes in docs/paper-gaps/.
% Every note loads this file via % Shared commands for the notes in docs/paper-gaps/.
% Every note loads this file via % Shared commands for the notes in docs/paper-gaps/.
% Every note loads this file via \input{command}. Project-wide notation
% belongs here, never in per-note redefinitions.

% --- Project configuration ------------------------------------------------
\newcommand{\projectIssueBase}{https://github.com/LionSR/QICLean/issues/}
\newcommand{\projectPrBase}{https://github.com/LionSR/QICLean/pull/}
\newcommand{\projectDocBase}{https://sirui-lu.com/QICLean/docs/}
% ---------------------------------------------------------------------------

\DeclareUnicodeCharacter{2080}{\ifmmode{}_{0}\else\textsubscript{0}\fi}
\DeclareUnicodeCharacter{2081}{\ifmmode{}_{1}\else\textsubscript{1}\fi}
\DeclareUnicodeCharacter{2082}{\ifmmode{}_{2}\else\textsubscript{2}\fi}
\DeclareUnicodeCharacter{2099}{\ifmmode{}_{n}\else\textsubscript{n}\fi}

\newcommand{\Lean}{\textsc{Lean}}

% A formal declaration name, upright sans, breakable at underscores.
\ExplSyntaxOn
\NewDocumentCommand{\leanid}{m}{
  \ifmmode
    \textsf{\detokenize{#1}}
  \else
    \group_begin:
    \urlstyle{sf}
    \tl_set:Nn \l_tmpa_tl {#1}
    \tl_replace_all:Nnn \l_tmpa_tl {\_} {_}
    \exp_args:NV \path \l_tmpa_tl
    \group_end:
  \fi
}
\ExplSyntaxOff

% Traceability links; use in footnotes, not in the running prose.
\newcommand{\ghissue}[1]{\href{\projectIssueBase#1}{GitHub issue \##1}}
\newcommand{\ghpr}[1]{\href{\projectPrBase#1}{PR \##1}}
\newcommand{\doclink}[1]{\href{\projectDocBase#1.html}{\path{#1}}}

% Domain notation, using \providecommand so a note-local refinement stays
% possible.
\providecommand{\tr}{\operatorname{tr}}

% Dirac notation and the identity operator, as in blueprint/src/macros.
\usepackage{dsfont}
\providecommand{\ket}[1]{|#1\rangle}
\providecommand{\bra}[1]{\langle #1|}
\providecommand{\braket}[2]{\langle #1 | #2 \rangle}
\providecommand{\ketbra}[2]{|#1\rangle\!\langle #2|}
\providecommand{\Id}{\mathds{1}}
\providecommand{\one}{\mathds{1}}
\providecommand{\C}{\mathbb{C}}
\providecommand{\MN}[1]{M_{#1}(\C)}
\providecommand{\MD}{M_D(\C)}

% Verdict marker: \gapnote{<kind>}{<status>}. Kind is the policy
% classification (clarification, local-correction, scope-restriction,
% unfaithful, false-source, open-gap); status is open, wip (elimination
% underway), resolved, or historical. Machine-read by the site index;
% prints a verdict line.
\newcommand{\gapnote}[2]{\par\noindent\textbf{Verdict:} #1 (#2).\par}
. Project-wide notation
% belongs here, never in per-note redefinitions.

% --- Project configuration ------------------------------------------------
\newcommand{\projectIssueBase}{https://github.com/LionSR/QICLean/issues/}
\newcommand{\projectPrBase}{https://github.com/LionSR/QICLean/pull/}
\newcommand{\projectDocBase}{https://sirui-lu.com/QICLean/docs/}
% ---------------------------------------------------------------------------

\DeclareUnicodeCharacter{2080}{\ifmmode{}_{0}\else\textsubscript{0}\fi}
\DeclareUnicodeCharacter{2081}{\ifmmode{}_{1}\else\textsubscript{1}\fi}
\DeclareUnicodeCharacter{2082}{\ifmmode{}_{2}\else\textsubscript{2}\fi}
\DeclareUnicodeCharacter{2099}{\ifmmode{}_{n}\else\textsubscript{n}\fi}

\newcommand{\Lean}{\textsc{Lean}}

% A formal declaration name, upright sans, breakable at underscores.
\ExplSyntaxOn
\NewDocumentCommand{\leanid}{m}{
  \ifmmode
    \textsf{\detokenize{#1}}
  \else
    \group_begin:
    \urlstyle{sf}
    \tl_set:Nn \l_tmpa_tl {#1}
    \tl_replace_all:Nnn \l_tmpa_tl {\_} {_}
    \exp_args:NV \path \l_tmpa_tl
    \group_end:
  \fi
}
\ExplSyntaxOff

% Traceability links; use in footnotes, not in the running prose.
\newcommand{\ghissue}[1]{\href{\projectIssueBase#1}{GitHub issue \##1}}
\newcommand{\ghpr}[1]{\href{\projectPrBase#1}{PR \##1}}
\newcommand{\doclink}[1]{\href{\projectDocBase#1.html}{\path{#1}}}

% Domain notation, using \providecommand so a note-local refinement stays
% possible.
\providecommand{\tr}{\operatorname{tr}}

% Dirac notation and the identity operator, as in blueprint/src/macros.
\usepackage{dsfont}
\providecommand{\ket}[1]{|#1\rangle}
\providecommand{\bra}[1]{\langle #1|}
\providecommand{\braket}[2]{\langle #1 | #2 \rangle}
\providecommand{\ketbra}[2]{|#1\rangle\!\langle #2|}
\providecommand{\Id}{\mathds{1}}
\providecommand{\one}{\mathds{1}}
\providecommand{\C}{\mathbb{C}}
\providecommand{\MN}[1]{M_{#1}(\C)}
\providecommand{\MD}{M_D(\C)}

% Verdict marker: \gapnote{<kind>}{<status>}. Kind is the policy
% classification (clarification, local-correction, scope-restriction,
% unfaithful, false-source, open-gap); status is open, wip (elimination
% underway), resolved, or historical. Machine-read by the site index;
% prints a verdict line.
\newcommand{\gapnote}[2]{\par\noindent\textbf{Verdict:} #1 (#2).\par}
. Project-wide notation
% belongs here, never in per-note redefinitions.

% --- Project configuration ------------------------------------------------
\newcommand{\projectIssueBase}{https://github.com/LionSR/QICLean/issues/}
\newcommand{\projectPrBase}{https://github.com/LionSR/QICLean/pull/}
\newcommand{\projectDocBase}{https://sirui-lu.com/QICLean/docs/}
% ---------------------------------------------------------------------------

\DeclareUnicodeCharacter{2080}{\ifmmode{}_{0}\else\textsubscript{0}\fi}
\DeclareUnicodeCharacter{2081}{\ifmmode{}_{1}\else\textsubscript{1}\fi}
\DeclareUnicodeCharacter{2082}{\ifmmode{}_{2}\else\textsubscript{2}\fi}
\DeclareUnicodeCharacter{2099}{\ifmmode{}_{n}\else\textsubscript{n}\fi}

\newcommand{\Lean}{\textsc{Lean}}

% A formal declaration name, upright sans, breakable at underscores.
\ExplSyntaxOn
\NewDocumentCommand{\leanid}{m}{
  \ifmmode
    \textsf{\detokenize{#1}}
  \else
    \group_begin:
    \urlstyle{sf}
    \tl_set:Nn \l_tmpa_tl {#1}
    \tl_replace_all:Nnn \l_tmpa_tl {\_} {_}
    \exp_args:NV \path \l_tmpa_tl
    \group_end:
  \fi
}
\ExplSyntaxOff

% Traceability links; use in footnotes, not in the running prose.
\newcommand{\ghissue}[1]{\href{\projectIssueBase#1}{GitHub issue \##1}}
\newcommand{\ghpr}[1]{\href{\projectPrBase#1}{PR \##1}}
\newcommand{\doclink}[1]{\href{\projectDocBase#1.html}{\path{#1}}}

% Domain notation, using \providecommand so a note-local refinement stays
% possible.
\providecommand{\tr}{\operatorname{tr}}

% Dirac notation and the identity operator, as in blueprint/src/macros.
\usepackage{dsfont}
\providecommand{\ket}[1]{|#1\rangle}
\providecommand{\bra}[1]{\langle #1|}
\providecommand{\braket}[2]{\langle #1 | #2 \rangle}
\providecommand{\ketbra}[2]{|#1\rangle\!\langle #2|}
\providecommand{\Id}{\mathds{1}}
\providecommand{\one}{\mathds{1}}
\providecommand{\C}{\mathbb{C}}
\providecommand{\MN}[1]{M_{#1}(\C)}
\providecommand{\MD}{M_D(\C)}

% Verdict marker: \gapnote{<kind>}{<status>}. Kind is the policy
% classification (clarification, local-correction, scope-restriction,
% unfaithful, false-source, open-gap); status is open, wip (elimination
% underway), resolved, or historical. Machine-read by the site index;
% prints a verdict line.
\newcommand{\gapnote}[2]{\par\noindent\textbf{Verdict:} #1 (#2).\par}

\title{The Scope of the Imaginary-Time Conjugation Bound for Interactions of Bounded Degree}
\author{}
\date{2026-11-19}
\begin{document}
\maketitle
\gapnote{scope-restriction}{open}

\section{The printed lemma}

Arad, Kuwahara and Landau,
\href{https://arxiv.org/abs/1406.3898}{arXiv:1406.3898v3}, consider a
Hamiltonian $H=\sum_Xh_X$ on finitely many particles with $h_X\ge0$,
$h_X=0$ for $|X|>k$, and the degree condition
\begin{align}
    \sum_{X\ni i}\lVert h_X\rVert\le g
    \label{eq:akl14_hadamard_degree}
\end{align}
at every particle $i$ (Section \texttt{sec:settings}, \texttt{low70.tex},
lines 400--445). For an operator $A$, let $E_A$ be a set of terms with
$[H,A]=\sum_{X\in E_A}[h_X,A]$ and $R=\sum_{X\in E_A}\lVert h_X\rVert$
(lines 523--525). Lemma \texttt{lem:Hadamard-bound} (lines 795--880) states
\begin{align}
    \lVert e^{sH}Ae^{-sH}\rVert\le\lVert A\rVert\,(1-gks)^{-R/(gk)}
    \qquad\text{for }0\le s<\frac1{gk}.
    \label{eq:akl14_hadamard_bound}
\end{align}
The proof expands $e^{sH}Ae^{-sH}$ in nested commutators, bounds
$\lVert[h_X,O]\rVert\le\lVert h_X\rVert\,\lVert O\rVert$ for $h_X\ge0$, and
bounds the norm sum of the terms entering the $j$-th commutator by
$R+(j-1)gk$.

\section{The formalized statement}

The formal statement\footnote{Formal declaration:
    \leanid{ManyBody.Interaction.norm_exp_conjugation_shift_le_two} in
    \doclink{QICLean/ManyBody/Dynamics/BoundedDegreeImaginaryTime}.}
assumes:
\begin{enumerate}
    \item every term satisfies $0\le h_X\le1$;
    \item every site lies in the supports of at most $g$ terms;
    \item the observable $O$ is supported on at most $k$ sites.
\end{enumerate}
It concludes, for every $z\in\C$ and $0\le s$ with $2gks\le1$,
\begin{align}
    \lVert e^{s(H-z)}Oe^{-s(H-z)}\rVert\le2\lVert O\rVert.
    \label{eq:akl14_hadamard_formal}
\end{align}
Under the first two assumptions, \eqref{eq:akl14_hadamard_degree} holds,
since every term meeting a site has norm at most one. Under the third,
the terms not commuting with $O$ meet one of at most $k$ sites, so
$R\le gk$. For $gks\le1/2$ this gives
$(1-gks)^{-R/(gk)}\le(1-gks)^{-1}\le2$, so
\eqref{eq:akl14_hadamard_formal} is the specialization of
\eqref{eq:akl14_hadamard_bound} to these hypotheses. The proof is the same
expansion: after $j$ commutators a piece is supported on at most $|S|+jk$
sites, at most $g(|S|+jk)$ terms fail to commute with it, and so
$\lVert\operatorname{ad}_H^m(O)\rVert\le(gk)^m\,m!\,\lVert O\rVert$ when
$|S|\le k$. The scalar shift $z$ does not change the commutators.

\section{What is not formalized}

The formal statement is narrower than \eqref{eq:akl14_hadamard_bound} in
three respects: the terms are contractions; the degree counts terms instead
of summing their norms, which is a stronger hypothesis when the norms are
small; and the observable lies on at most $k$ sites, with only the
consequence $2\lVert O\rVert$ at $s\le1/(2gk)$ stated. The printed bound for
an arbitrary $R$ and $s<1/(gk)$, with the norm-weighted degree
\eqref{eq:akl14_hadamard_degree}, remains unformalized. Its elimination is a
formal statement with $R$ and the norm-weighted degree as hypotheses, from
the bound $\lVert\operatorname{ad}_H^m(A)\rVert\le\prod_{j<m}(R+jgk)\,
\lVert A\rVert$ and the binomial series of $(1-gks)^{-R/(gk)}$. The hard
truncation of a region\footnote{Formal declaration:
    \leanid{ManyBody.Interaction.exists_region_hard_truncation_spectral_comparison}
    in \doclink{QICLean/ManyBody/Hamiltonian/HardTruncation}.}
uses \eqref{eq:akl14_hadamard_formal} only for terms supported on at most
$k$ sites, where the restricted form suffices.
\end{document}
